\section{Giới thiệu}
\label{sec:intro}

Công nghệ container hóa đã trở thành nền tảng của điện toán đám mây hiện đại nhờ tính hiệu quả và khả năng mở rộng. Không giống như các máy ảo truyền thống, container chia sẻ nhân hệ điều hành (OS kernel) của máy chủ, điều này làm giảm chi phí vận hành nhưng cũng đưa ra các rủi ro bảo mật nghiêm trọng thông qua lời gọi hệ thống (syscall)---giao diện giữa không gian người dùng và nhân hệ điều hành~\cite{ref1_jarkas_survey,ref3_sultan_survey}.

Syscall mở ra một bề mặt tấn công rộng lớn, khiến chúng trở thành mục tiêu hấp dẫn cho các hành vi leo thang đặc quyền, thoát container và nhiều khai thác khác~\cite{ref5_combe_docker}. Các giải pháp hiện có như hồ sơ seccomp chỉ cung cấp mức độ bảo vệ hạn chế. Các hồ sơ tĩnh khó duy trì và có thể chặn các thao tác hợp lệ hoặc không phát hiện được các mối đe dọa tinh vi.

Để khắc phục những hạn chế này, chúng tôi đề xuất DeSFAM, Lọc lời gọi hệ thống động dựa trên eBPF và giảm thiểu bất thường. DeSFAM cung cấp lập hồ sơ syscall thời gian thực, phát hiện bất thường sâu và thực thi chính sách thích nghi. Framework này tích hợp:
\begin{enumerate}[(i)]
  \item Lập danh sách truy cập syscall lai sử dụng phân tích tĩnh và phân tích động dựa trên eBPF;
  \item SyscallAD, một công cụ phát hiện kết hợp Bộ mã hóa biến phân tự động (VAE) và Rừng cô lập
        (iForest) để xác định các chuỗi syscall bất thường; và
  \item Quản lý chính sách thích nghi sử dụng ánh xạ khung chiến thuật, kỹ thuật và kiến thức đối nghịch
        MITRE (ATT\&CK), tương quan lỗ hổng và phơi nhiễm chung (CVE) và các móc eBPF LSM để kiểm
        soát syscall theo ngữ cảnh.
\end{enumerate}

DeSFAM thực thi kiểm soát syscall chi tiết bằng cách sử dụng bản đồ eBPF và các móc LSM, chặn động các syscall có nguy cơ cao (ví dụ: \texttt{execve}, \texttt{ptrace}) và cập nhật chính sách tại thời điểm chạy mà không cần khởi động lại container. Framework này kết hợp một quy trình lập hồ sơ syscall lai với SyscallAD, một bộ phát hiện bất thường có độ trễ thấp dựa trên Bộ mã hóa biến phân tự động (VAE) và Rừng cô lập. Để ưu tiên các mối đe dọa, DeSFAM tích hợp tính điểm rủi ro theo ngữ cảnh bằng cách tương quan hành vi syscall với các chiến thuật MITRE ATT\&CK và các CVE đã biết. Đánh giá sử dụng bộ dữ liệu DongTing và các kịch bản CVE thực tế cho thấy độ chính xác 94\%, độ thu hồi 90\%, độ trễ thực thi dưới mili-giây và chi phí hiệu năng dưới 1\%.

Các đóng góp chính của chúng tôi bao gồm:
\begin{itemize}
  \item Một quy trình lập hồ sơ lai để tạo danh sách cho phép syscall.
  \item SyscallAD, một bộ phát hiện bất thường VAE + iForest có độ trễ thấp.
  \item Tính điểm rủi ro theo ngữ cảnh sử dụng MITRE ATT\&CK và tương quan CVE.
  \item Thực thi syscall chi tiết sử dụng bản đồ eBPF và các móc LSM.
  \item Vòng phản hồi thời gian chạy để tự động tăng cường chính sách dựa trên các vi phạm.
\end{itemize}

Bài báo được tổ chức như sau: Mục~\ref{sec:background} trình bày các kiến thức nền tảng về lọc syscall và eBPF. Mục~\ref{sec:related} thảo luận các công trình liên quan. Mục~\ref{sec:desfam} mô tả kiến trúc của DeSFAM. Mục~\ref{sec:evaluation} trình bày kết quả đánh giá, tiếp theo là thảo luận ở Mục~\ref{sec:discussion} và kết luận ở Mục~\ref{sec:conclusion}.
